\documentclass[letterpaper]{article} % DO NOT CHANGE THIS
\usepackage{aaai2027} % DO NOT CHANGE THIS
\usepackage[hyphens]{url} % DO NOT CHANGE THIS
\usepackage{graphicx} % DO NOT CHANGE THIS
\urlstyle{rm} % DO NOT CHANGE THIS
\def\UrlFont{\rm} % DO NOT CHANGE THIS
\usepackage{natbib} % DO NOT CHANGE THIS
\usepackage{caption} % DO NOT CHANGE THIS
\frenchspacing % DO NOT CHANGE THIS
\usepackage{amsmath}
\usepackage{amssymb}
\usepackage{booktabs}
\usepackage{graphicx}
\usepackage[hyphens]{url}
\usepackage[ruled,vlined,linesnumbered,portuguese]{algorithm2e}
\usepackage{microtype}

\pdfinfo{
/TemplateVersion (2027.1)
}

% Optional: Set counter for section numbers if needed (0 means no numbers, 1 means section numbers, 2 means sections and subsections)
\setcounter{secnumdepth}{2}

% Custom commands for diacritics / accented characters
\newcommand{\aacute}{\'a}
\newcommand{\eacute}{\'e}
\newcommand{\iacute}{\'i}
\newcommand{\oacute}{\'o}
\newcommand{\uacute}{\'u}
\newcommand{\agrave}{\`a}
\newcommand{\acirc}{\^a}
\newcommand{\ecirc}{\^e}
\newcommand{\ocirc}{\^o}
\newcommand{\atilde}{\~a}
\newcommand{\otilde}{\~o}
\newcommand{\ccedil}{\c{c}}
\newcommand{\Aacute}{\'A}
\newcommand{\Eacute}{\'E}
\newcommand{\Iacute}{\'I}
\newcommand{\Oacute}{\'O}
\newcommand{\Uacute}{\'U}
\newcommand{\Agrave}{\`A}
\newcommand{\Acirc}{\^A}
\newcommand{\Ecirc}{\^E}
\newcommand{\Ocirc}{\^O}
\newcommand{\Atilde}{\~A}
\newcommand{\Otilde}{\~O}
\newcommand{\Ccedil}{\c{C}}

\renewcommand{\a}{{\agrave}}
\newcommand{\abstracao}{abstra{\c{c}}{\~a}o}
\newcommand{\aceleracao}{acelera{\ccedil}{\atilde}o}
\newcommand{\acrescimo}{acr{\'{e}}scimo}
\newcommand{\Acuracia}{Acur{\aacute}cia}
\newcommand{\acuracia}{acur{\aacute}cia}
\newcommand{\alcancando}{alcan{\ccedil}ando}
\newcommand{\alcancou}{alcan{\c{c}}ou}
\newcommand{\Alem}{Al{\eacute}m}
\newcommand{\alem}{al{\eacute}m}
\newcommand{\alteracao}{altera{\c{c}}{\~a}o}
\newcommand{\alucinacoes}{alucina{\ccedil}{\otilde}es}
\newcommand{\Analise}{An{\aacute}lise}
\newcommand{\analise}{an{\aacute}lise}
\newcommand{\analiticas}{anal{\iacute}ticas}
\newcommand{\analitico}{anal{\iacute}tico}
\newcommand{\ancora}{{\^{a}}ncora}
\newcommand{\AntiTrapaca}{Anti-Trapa{\ccedil}a}
\newcommand{\antitrapaca}{anti-trapa{\ccedil}a}
\newcommand{\aplicacao}{aplica{\ccedil}{\atilde}o}
\newcommand{\apos}{ap{\oacute}s}
\newcommand{\aproximacao}{aproxima{\ccedil}{\atilde}o}
\newcommand{\aquisicao}{aquisi{\ccedil}{\atilde}o}
\newcommand{\arcabouco}{arcabou{\ccedil}o}
\newcommand{\as}{{\agrave}s}
\newcommand{\assintotica}{assint{\oacute}tica}
\newcommand{\ate}{at{\eacute}}
\newcommand{\atribuicao}{atribui{\ccedil}{\atilde}o}
\newcommand{\ausencia}{aus{\^{e}}ncia}
\newcommand{\automatico}{autom{\aacute}tico}
\newcommand{\Avaliacao}{Avalia{\ccedil}{\atilde}o}
\newcommand{\avaliacao}{avalia{\ccedil}{\atilde}o}
\newcommand{\avaliacoes}{avalia{\ccedil}{\otilde}es}
\newcommand{\avancada}{avan{\c{c}}ada}
\newcommand{\avancado}{avan{\ccedil}ado}
\newcommand{\avancados}{avan{\ccedil}ados}
\newcommand{\avanco}{avan{\c{c}}o}
\newcommand{\avancos}{avan{\c{c}}os}
\newcommand{\avancou}{avan{\c{c}}ou}
\newcommand{\balanco}{balan{\ccedil}o}
\newcommand{\basica}{b{\aacute}sica}
\newcommand{\beneficios}{benef{\'{i}}cios}
\newcommand{\binario}{bin{\'{a}}rio}
\newcommand{\carencia}{car{\^{e}}ncia}
\newcommand{\catastrofica}{catastr{\'{o}}fica}
\newcommand{\cenarios}{cen{\aacute}rios}
\newcommand{\centralizacao}{centraliza{\ccedil}{\atilde}o}
\newcommand{\cientifica}{cient{\iacute}fica}
\newcommand{\classificacao}{classifica{\ccedil}{\atilde}o}
\newcommand{\Codigo}{C{\oacute}digo}
\newcommand{\codigo}{c{\oacute}digo}
\newcommand{\codigos}{c{\oacute}digos}
\newcommand{\coerencia}{coer{\ecirc}ncia}
\newcommand{\combinacao}{combina{\c{c}}{\~a}o}
\newcommand{\comparacao}{compara{\ccedil}{\atilde}o}
\newcommand{\compativeis}{compat{\iacute}veis}
\newcommand{\composicao}{composi{\ccedil}{\atilde}o}
\newcommand{\concluida}{conclu{\iacute}da}
\newcommand{\Conclusao}{Conclus{\atilde}o}
\newcommand{\conclusao}{conclus{\~a}o}
\newcommand{\conclusoes}{conclus{\~o}es}
\newcommand{\Conclusoes}{Conclus{\~o}es}
\newcommand{\condicoes}{condi{\ccedil}{\otilde}es}
\newcommand{\confianca}{confian{\ccedil}a}
\newcommand{\constroi}{constr{\oacute}i}
\newcommand{\contencao}{conten{\ccedil}{\atilde}o}
\newcommand{\Contingencia}{Conting{\ecirc}ncia}
\newcommand{\contingencia}{conting{\ecirc}ncia}
\newcommand{\continuo}{cont{\iacute}nuo}
\newcommand{\contrario}{contr{\aacute}rio}
\newcommand{\Convergencia}{Converg{\ecirc}ncia}
\newcommand{\convergencia}{converg{\ecirc}ncia}
\newcommand{\correcao}{corre{\ccedil}{\atilde}o}
\newcommand{\Criterio}{Crit{\eacute}rio}
\newcommand{\criterio}{crit{\eacute}rio}
\newcommand{\decompoe}{decomp{\otilde}e}
\newcommand{\Decomposicao}{Decomposi{\ccedil}{\atilde}o}
\newcommand{\decomposicao}{decomposi{\ccedil}{\atilde}o}
\newcommand{\decorrencia}{decorr{\ecirc}ncia}
\newcommand{\deducao}{dedu{\ccedil}{\atilde}o}
\newcommand{\definicao}{defini{\ccedil}{\atilde}o}
\newcommand{\delimitacao}{delimita{\ccedil}{\atilde}o}
\newcommand{\demonstracao}{demonstra{\ccedil}{\atilde}o}
\newcommand{\demonstracoes}{demonstra{\ccedil}{\otilde}es}
\newcommand{\dependencias}{depend{\^{e}}ncias}
\newcommand{\descalibracoes}{descalibra{\c{c}}{\~o}es}
\newcommand{\descricao}{descri{\ccedil}{\atilde}o}
\newcommand{\deteccao}{detec{\ccedil}{\atilde}o}
\newcommand{\deterministico}{determin{\iacute}stico}
\newcommand{\deterministicos}{determin{\'{i}}sticos}
\newcommand{\dialogo}{di{\aacute}logo}
\newcommand{\diaria}{di{\aacute}ria}
\newcommand{\diarias}{di{\aacute}rias}
\newcommand{\Diferenca}{Diferen{\ccedil}a}
\newcommand{\diferenca}{diferen{\ccedil}a}
\newcommand{\diferencas}{diferen{\ccedil}as}
\newcommand{\dinamica}{din{\acirc}mica}
\newcommand{\dinamicas}{din{\acirc}micas}
\newcommand{\direcao}{dire{\c{c}}{\~a}o}
\newcommand{\discrepancia}{discrep{\^{a}}ncia}
\newcommand{\discrepancias}{discrep{\^{a}}ncias}
\newcommand{\Discussao}{Discuss{\atilde}o}
\newcommand{\discussao}{discuss{\~a}o}
\newcommand{\dispendio}{disp{\ecirc}ndio}
\newcommand{\disponibilizacao}{disponibiliza{\ccedil}{\atilde}o}
\newcommand{\distribuicao}{distribui{\ccedil}{\atilde}o}
\newcommand{\distribuicoes}{distribui{\ccedil}{\otilde}es}
\newcommand{\distribuidos}{distribu{\'{i}}dos}
\newcommand{\divergencia}{diverg{\ecirc}ncia}
\newcommand{\divergencias}{diverg{\ecirc}ncias}
\newcommand{\dominio}{dom{\iacute}nio}
\newcommand{\e}{{\eacute}}
\newcommand{\eficacia}{efic{\aacute}cia}
\newcommand{\Eficacia}{Efic{\aacute}cia}
\newcommand{\eficiencia}{efici{\ecirc}ncia}
\newcommand{\elevacao}{eleva{\ccedil}{\atilde}o}
\newcommand{\Elicitacao}{Elicita{\ccedil}{\atilde}o}
\newcommand{\elicitacao}{elicita{\ccedil}{\atilde}o}
\newcommand{\emissao}{emiss{\atilde}o}
\newcommand{\empirica}{emp{\iacute}rica}
\newcommand{\Empirico}{Emp{\iacute}rico}
\newcommand{\empirico}{emp{\iacute}rico}
\newcommand{\empiricos}{emp{\iacute}ricos}
\newcommand{\entao}{ent{\~a}o}
\newcommand{\equilibrio}{equil{\iacute}brio}
\newcommand{\equivalencia}{equival{\ecirc}ncia}
\newcommand{\escalavel}{escal{\'{a}}vel}
\newcommand{\esforco}{esfor{\ccedil}o}
\newcommand{\espaco}{espa{\ccedil}o}
\newcommand{\espacos}{espa{\c{c}}os}
\newcommand{\especificacao}{especifica{\ccedil}{\atilde}o}
\newcommand{\especificas}{espec{\iacute}ficas}
\newcommand{\espuria}{esp{\'{u}}ria}
\newcommand{\espurias}{esp{\uacute}rias}
\newcommand{\Espurio}{Esp{\uacute}rio}
\newcommand{\espurio}{esp{\uacute}rio}
\newcommand{\Estatistica}{Estat{\iacute}stica}
\newcommand{\estatistica}{estat{\iacute}stica}
\newcommand{\estatistico}{estat{\iacute}stico}
\newcommand{\estatisticos}{estat{\iacute}sticos}
\newcommand{\estavel}{est{\'{a}}vel}
\newcommand{\estimacao}{estima{\ccedil}{\atilde}o}
\newcommand{\estimulo}{est{\'{i}}mulo}
\newcommand{\estocastica}{estoc{\aacute}stica}
\newcommand{\estocasticas}{estoc{\'{a}}sticas}
\newcommand{\estocastico}{estoc{\aacute}stico}
\newcommand{\estrategia}{estrat{\'{e}}gia}
\newcommand{\estrategias}{estrat{\eacute}gias}
\newcommand{\evolucao}{evolu{\ccedil}{\atilde}o}
\newcommand{\Execucao}{Execu{\ccedil}{\atilde}o}
\newcommand{\execucao}{execu{\ccedil}{\atilde}o}
\newcommand{\execucoes}{execu{\ccedil}{\otilde}es}
\newcommand{\executaveis}{execut{\'{a}}veis}
\newcommand{\executavel}{execut{\aacute}vel}
\newcommand{\existencia}{exist{\^{e}}ncia}
\newcommand{\explicita}{expl{\iacute}cita}
\newcommand{\explicitacao}{explicita{\ccedil}{\atilde}o}
\newcommand{\explicitas}{expl{\iacute}citas}
\newcommand{\expoe}{exp{\otilde}e}
\newcommand{\expressao}{express{\~a}o}
\newcommand{\extensoes}{extens{\otilde}es}
\newcommand{\familias}{fam{\iacute}lias}
\newcommand{\fenomeno}{fen{\ocirc}meno}
\newcommand{\fenomenos}{fen{\ocirc}menos}
\newcommand{\finalizacao}{finaliza{\ccedil}{\atilde}o}
\newcommand{\flutuacoes}{flutua{\c{c}}{\~o}es}
\newcommand{\forca}{for{\ccedil}a}
\newcommand{\forcar}{for{\ccedil}ar}
\newcommand{\Formulacao}{Formula{\ccedil}{\atilde}o}
\newcommand{\formulacao}{formula{\ccedil}{\atilde}o}
\newcommand{\frequencia}{frequ{\ecirc}ncia}
\newcommand{\funcao}{fun{\ccedil}{\atilde}o}
\newcommand{\funcoes}{fun{\ccedil}{\otilde}es}
\newcommand{\Fundamentacao}{Fundamenta{\ccedil}{\atilde}o}
\newcommand{\fundamentacao}{fundamenta{\ccedil}{\atilde}o}
\newcommand{\generalizacao}{generaliza{\ccedil}{\atilde}o}
\newcommand{\geometrica}{geom{\'{e}}trica}
\newcommand{\geometricas}{geom{\eacute}tricas}
\newcommand{\geometricosemantica}{geom{\eacute}trico-sem{\acirc}ntica}
\newcommand{\geracao}{gera{\ccedil}{\atilde}o}
\newcommand{\Gestao}{Gest{\atilde}o}
\newcommand{\gestao}{gest{\atilde}o}
\newcommand{\hegemonico}{hegem{\^{o}}nico}
\newcommand{\hermetico}{herm{\'{e}}tico}
\newcommand{\heuristicos}{heur{\'{i}}sticos}
\newcommand{\Hipotese}{Hip{\oacute}tese}
\newcommand{\hipotese}{hip{\oacute}tese}
\newcommand{\Hipoteses}{Hip{\oacute}teses}
\newcommand{\hipoteses}{hip{\oacute}teses}
\newcommand{\identicas}{id{\ecirc}nticas}
\newcommand{\impar}{{\'{i}}mpar}
\newcommand{\impoe}{imp{\otilde}e}
\newcommand{\importacao}{importa{\ccedil}{\atilde}o}
\newcommand{\importacoes}{importa{\ccedil}{\otilde}es}
\newcommand{\incidencia}{incid{\^{e}}ncia}
\newcommand{\incongruencias}{incongru{\^{e}}ncias}
\newcommand{\incorrecoes}{incorre{\ccedil}{\otilde}es}
\newcommand{\indexacao}{indexa{\c{c}}{\~a}o}
\newcommand{\Indice}{{\Iacute}ndice}
\newcommand{\indice}{{\iacute}ndice}
\newcommand{\indices}{{\iacute}ndices}
\newcommand{\indispensavel}{indispens{\'{a}}vel}
\newcommand{\ineditas}{in{\'{e}}ditas}
\newcommand{\Inferencia}{Infer{\ecirc}ncia}
\newcommand{\inferencia}{infer{\ecirc}ncia}
\newcommand{\insercao}{inser{\ccedil}{\atilde}o}
\newcommand{\instrucoes}{instru{\ccedil}{\otilde}es}
\newcommand{\instruidos}{instru{\iacute}dos}
\newcommand{\integracao}{integra{\ccedil}{\atilde}o}
\newcommand{\inteligencia}{intelig{\ecirc}ncia}
\newcommand{\interacao}{intera{\c{c}}{\~a}o}
\newcommand{\Intermediaria}{Intermedi{\'{a}}ria}
\newcommand{\intermediarias}{intermedi{\'{a}}rias}
\newcommand{\intermediario}{intermedi{\'{a}}rio}
\newcommand{\intermediarios}{intermedi{\aacute}rios}
\newcommand{\interrupcoes}{interrup{\ccedil}{\otilde}es}
\newcommand{\Introducao}{Introdu{\ccedil}{\atilde}o}
\newcommand{\introducao}{introdu{\ccedil}{\atilde}o}
\newcommand{\intuia}{intu{\'{i}}a}
\newcommand{\invariancia}{invari{\acirc}ncia}
\newcommand{\inversao}{invers{\~a}o}
\newcommand{\investigacao}{investiga{\c{c}}{\~a}o}
\newcommand{\iteracao}{itera{\ccedil}{\atilde}o}
\newcommand{\iteracoes}{itera{\ccedil}{\otilde}es}
\newcommand{\ja}{j{\aacute}}
\newcommand{\Ja}{J{\'{a}}}
\newcommand{\Latencia}{Lat{\^{e}}ncia}
\newcommand{\latencia}{lat{\^{e}}ncia}
\newcommand{\logica}{l{\oacute}gica}
\newcommand{\logicogeometricas}{l{\'{o}}gico-geom{\'{e}}tricas}
\newcommand{\mandatoria}{mandat{\'{a}}ria}
\newcommand{\mandatorias}{mandat{\'{a}}rias}
\newcommand{\mandatorio}{mandat{\oacute}rio}
\newcommand{\maximo}{m{\aacute}ximo}
\newcommand{\medias}{m{\'{e}}dias}
\newcommand{\medio}{m{\eacute}dio}
\newcommand{\memorizacao}{memoriza{\ccedil}{\atilde}o}
\newcommand{\mensuracao}{mensura{\ccedil}{\atilde}o}
\newcommand{\metodo}{m{\eacute}todo}
\newcommand{\metodologica}{metodol{\oacute}gica}
\newcommand{\metodologico}{metodol{\oacute}gico}
\newcommand{\metodos}{m{\eacute}todos}
\newcommand{\metricas}{m{\eacute}tricas}
\newcommand{\microscopica}{microsc{\'{o}}pica}
\newcommand{\minimo}{m{\iacute}nimo}
\newcommand{\modulos}{m{\oacute}dulos}
\newcommand{\multiplas}{m{\uacute}ltiplas}
\newcommand{\nao}{n{\atilde}o}
\newcommand{\Nao}{N{\atilde}o}
\newcommand{\naoparametrica}{n{\~a}o-param{\'{e}}trica}
\newcommand{\naoparametrico}{n{\atilde}o-param{\eacute}trico}
\newcommand{\naoparametricos}{n{\atilde}o-param{\eacute}tricos}
\newcommand{\naosuperioridade}{n{\~a}o-superioridade}
\newcommand{\necessaria}{necess{\aacute}ria}
\newcommand{\necessario}{necess{\aacute}rio}
\newcommand{\niveis}{n{\'{i}}veis}
\newcommand{\nivel}{n{\iacute}vel}
\newcommand{\notavel}{not{\aacute}vel}
\newcommand{\numero}{n{\uacute}mero}
\newcommand{\observacoes}{observa{\ccedil}{\otilde}es}
\newcommand{\Obtencao}{Obten{\ccedil}{\atilde}o}
\newcommand{\obtencao}{obten{\ccedil}{\atilde}o}
\newcommand{\ocorrencia}{ocorr{\ecirc}ncia}
\newcommand{\operacoes}{opera{\ccedil}{\otilde}es}
\newcommand{\orcamento}{or{\c{c}}amento}
\newcommand{\oscilacoes}{oscila{\ccedil}{\otilde}es}
\newcommand{\otica}{{\'{o}}tica}
\newcommand{\padrao}{padr{\atilde}o}
\newcommand{\padroes}{padr{\otilde}es}
\newcommand{\padronizacao}{padroniza{\c{c}}{\~a}o}
\newcommand{\parametros}{par{\acirc}metros}
\newcommand{\passiveis}{pass{\iacute}veis}
\newcommand{\permutacao}{permuta{\ccedil}{\atilde}o}
\newcommand{\Persistencia}{Persist{\ecirc}ncia}
\newcommand{\porem}{por{\eacute}m}
\newcommand{\pratica}{pr{\'{a}}tica}
\newcommand{\presenca}{presen{\ccedil}a}
\newcommand{\previa}{pr{\'{e}}via}
\newcommand{\previas}{pr{\'{e}}vias}
\newcommand{\previo}{pr{\'{e}}vio}
\newcommand{\primario}{prim{\'{a}}rio}
\newcommand{\programacao}{programa{\ccedil}{\atilde}o}
\newcommand{\Proibicao}{Proibi{\ccedil}{\atilde}o}
\newcommand{\propensao}{propens{\~a}o}
\newcommand{\proprietarios}{propriet{\aacute}rios}
\newcommand{\quantificacao}{quantifica{\c{c}}{\~a}o}
\newcommand{\raciocinio}{racioc{\iacute}nio}
\newcommand{\ramificacoes}{ramifica{\ccedil}{\otilde}es}
\newcommand{\realimentacao}{realimenta{\ccedil}{\atilde}o}
\newcommand{\recemapresentado}{rec{\'{e}}m-apresentado}
\newcommand{\Recuperacao}{Recupera{\ccedil}{\atilde}o}
\newcommand{\recuperacao}{recupera{\ccedil}{\atilde}o}
\newcommand{\recuperacoes}{recupera{\ccedil}{\otilde}es}
\newcommand{\reducao}{redu{\ccedil}{\atilde}o}
\newcommand{\referencia}{refer{\ecirc}ncia}
\newcommand{\reflexao}{reflex{\~a}o}
\newcommand{\Regressao}{Regress{\atilde}o}
\newcommand{\regressao}{regress{\atilde}o}
\newcommand{\regressoes}{regress{\otilde}es}
\newcommand{\relacao}{rela{\ccedil}{\atilde}o}
\newcommand{\relevancia}{relev{\^{a}}ncia}
\newcommand{\repertorio}{repert{\oacute}rio}
\newcommand{\replicas}{r{\eacute}plicas}
\newcommand{\reposicao}{reposi{\ccedil}{\atilde}o}
\newcommand{\Repositorio}{Reposit{\oacute}rio}
\newcommand{\Representacao}{Representa{\ccedil}{\atilde}o}
\newcommand{\representacao}{representa{\ccedil}{\atilde}o}
\newcommand{\reprodutiveis}{reprodut{\iacute}veis}
\newcommand{\reprodutivel}{reprodut{\'{i}}vel}
\newcommand{\reproduzivel}{reproduz{\iacute}vel}
\newcommand{\requisicao}{requisi{\ccedil}{\atilde}o}
\newcommand{\requisicoes}{requisi{\ccedil}{\otilde}es}
\newcommand{\resiliencia}{resili{\ecirc}ncia}
\newcommand{\resolucao}{resolu{\ccedil}{\atilde}o}
\newcommand{\restricao}{restri{\ccedil}{\atilde}o}
\newcommand{\restricoes}{restri{\ccedil}{\otilde}es}
\newcommand{\retroalimentacao}{retroalimenta{\c{c}}{\~a}o}
\newcommand{\rigidas}{r{\iacute}gidas}
\newcommand{\rigido}{r{\iacute}gido}
\newcommand{\ruido}{ru{\iacute}do}
\newcommand{\saida}{sa{\'{i}}da}
\newcommand{\saidas}{sa{\iacute}das}
\newcommand{\sao}{s{\atilde}o}
\newcommand{\saturacao}{satura{\ccedil}{\atilde}o}
\newcommand{\Secao}{Se{\ccedil}{\atilde}o}
\newcommand{\secao}{se{\c{c}}{\~a}o}
\newcommand{\secoes}{se{\c{c}}{\~o}es}
\newcommand{\Secoes}{Se{\c{c}}{\~o}es}
\newcommand{\selecao}{sele{\c{c}}{\~a}o}
\newcommand{\Semantica}{Sem{\acirc}ntica}
\newcommand{\semantica}{sem{\acirc}ntica}
\newcommand{\semanticas}{sem{\acirc}nticas}
\newcommand{\Semantico}{Sem{\acirc}ntico}
\newcommand{\semantico}{sem{\acirc}ntico}
\newcommand{\semanticos}{sem{\acirc}nticos}
\newcommand{\sensivel}{sens{\iacute}vel}
\newcommand{\significancia}{signific{\acirc}ncia}
\newcommand{\simbolica}{simb{\'{o}}lica}
\newcommand{\simultanea}{simult{\acirc}nea}
\newcommand{\simultaneo}{simult{\^{a}}neo}
\newcommand{\sintatica}{sint{\'{a}}tica}
\newcommand{\sintaticas}{sint{\aacute}ticas}
\newcommand{\Sintese}{S{\iacute}ntese}
\newcommand{\sintese}{s{\iacute}ntese}
\newcommand{\sintetico}{sint{\eacute}tico}
\newcommand{\sistematicas}{sistem{\aacute}ticas}
\newcommand{\situacao}{situa{\ccedil}{\atilde}o}
\newcommand{\solucao}{solu{\ccedil}{\atilde}o}
\newcommand{\solucoes}{solu{\ccedil}{\otilde}es}
\newcommand{\tambem}{tamb{\eacute}m}
\newcommand{\teoricoestatistica}{te{\oacute}rico-estat{\iacute}stica}
\newcommand{\tolerancia}{toler{\acirc}ncia}
\newcommand{\transformacao}{transforma{\ccedil}{\atilde}o}
\newcommand{\transformacoes}{transforma{\c{c}}{\~o}es}
\newcommand{\Transicoes}{Transi{\ccedil}{\otilde}es}
\newcommand{\transicoes}{transi{\ccedil}{\otilde}es}
\newcommand{\translacao}{transla{\ccedil}{\atilde}o}
\newcommand{\tres}{tr{\^{e}}s}
\newcommand{\unica}{{\uacute}nica}
\newcommand{\unico}{{\'{u}}nico}
\newcommand{\util}{{\'{u}}til}
\newcommand{\validacao}{valida{\ccedil}{\atilde}o}
\newcommand{\Validacao}{Valida{\ccedil}{\atilde}o}
\newcommand{\validas}{v{\'{a}}lidas}
\newcommand{\valido}{v{\'{a}}lido}
\newcommand{\variacao}{varia{\ccedil}{\atilde}o}
\newcommand{\Variacoes}{Varia{\ccedil}{\otilde}es}
\newcommand{\variacoes}{varia{\ccedil}{\otilde}es}
\newcommand{\variavel}{vari{\aacute}vel}
\newcommand{\vazao}{vaz{\atilde}o}
\newcommand{\verificacao}{verifica{\ccedil}{\atilde}o}
\newcommand{\visiveis}{vis{\iacute}veis}
\newcommand{\visivel}{vis{\iacute}vel}


\title{\textbf{S\'{\i}ntese de Programas com CEGIS e Restri\c{c}\~{o}es Anti-Overfitting no ARC-AGI: Uma \Analise Estrutural e Estat\'{\i}stica}}

\author{
    Fabio Luiz da Costa Cieslak\equalcontrib\textsuperscript{1},
    Izadora Candotti de Oliveira\equalcontrib\textsuperscript{1},
    Matheus Machado\equalcontrib\textsuperscript{1}
}
\affiliations{
    Universidade Federal do Rio Grande do Sul\textsuperscript{1}
}

\begin{document}

\maketitle

% Seção 0: Resumo
\begin{abstract}
Resolver problemas visuais e l\^{o}gicos do benchmark ARC-AGI por meio da gera\c{c}\~{a}o autom\'{a}tica de c\'{o}digo em Python usando modelos de linguagem (LLMs) \'{e} um desafio central para a intelig\^{e}ncia artificial. Um problema recorrente nessa abordagem \'{e} o ajuste excessivo aos exemplos de treino (\textit{overfitting}): o modelo cria solu\c{c}\~{o}es que funcionam apenas para os exemplos mostrados, mas falham no teste. Neste trabalho, investigamos o m\'{e}todo de S\'{\i}ntese Indutiva Guiada por Contraexemplos (CEGIS) e avaliamos uma varia\c{c}\~{a}o com regras anti-ajuste (\textit{AntiCheat}), que pro\'{\i}be regras pontuais para pixels espec\'{\i}ficos e exige que o modelo explique em texto a l\'{o}gica geral antes de escrever o c\'{o}digo. Avaliamos 8 modelos de linguagem em 100 tarefas desafiadoras do ARC-AGI. Os resultados mostram que a inclus\~{a}o das regras anti-ajuste n\~{a}o alterou a acur\'{a}cia global de forma estatisticamente significante (acur\'{a}cia de 50,4\% no CEGIS contra 49,8\% no AntiCheat, com $p = 0,5515$). A an\'{a}lise qualitativa do c\'{o}digo gerado tamb\'{e}m indicou que ambos os m\'{e}todos produzem programas estruturalmente semelhantes e falham pelos mesmos motivos, principalmente pela incapacidade de encontrar a l\'{o}gica correta dentro do limite de tentativas. Por outro lado, a an\'{a}lise estrutural revelou que as tarefas seguem uma escala de dificuldade unidimensional quase perfeita de Guttman ($CR = 0,9738$) e que os c\'{o}digos com sobreajuste violam a Navalha de Occam, sendo 1,61 vezes maiores do que os programas que acertam de primeira ($p = 5,15 \times 10^{-9}$). Por fim, observamos que 86,0\% das corre\c{c}\~{o}es bem-sucedidas do CEGIS ocorrem at\'{e} a terceira itera\c{c}\~{a}o.
\end{abstract}



% Seção 1: Introdução
\section{Introdu\c{c}\~{a}o}

O \textit{Abstraction and Reasoning Corpus} (ARC-AGI) \cite{chollet2019measure} foi criado para medir a capacidade de sistemas artificiais de aprender novas regras a partir de poucos exemplos. Cada tarefa do ARC apresenta grades bidimensionais com n\'{u}meros que representam cores e desafia o sistema a descobrir a regra de transforma\c{c}\~{a}o para gerar a grade de sa\'{\i}da a partir de uma grade de entrada in\'{e}dita.

Uma das abordagens mais promissoras para resolver o ARC \'{e} a s\'{\i}ntese de programas com Modelos de Linguagem (LLMs) \cite{greenblatt2024getting}. Em vez de gerar a grade diretamente como texto, o modelo escreve uma fun\c{c}\~{a}o em Python (\texttt{transform(grid)}) que executa a l\'{o}gica da tarefa. Para melhorar o c\'{o}digo gerado, \'{e} comum utilizar a S\'{\i}ntese Indutiva Guiada por Contraexemplos (CEGIS) \cite{solar2008program}: o c\'{o}digo \'{e} testado nos exemplos de treino; se falhar em algum exemplo (chamado de contraexemplo), o erro \'{e} devolvido ao modelo para que ele fa\c{c}a uma corre\c{c}\~{a}o iterativa.

No entanto, essa abordagem enfrenta um desafio cr\'{\i}tico: o ajuste excessivo aos dados de treino (\textit{overfitting}). O modelo pode consertar o c\'{o}digo inserindo regras pontuais (\textit{hardcoding}) ou condicionais especiais para os exemplos de treino conhecidos, sem aprender o padr\~{a}o l\'{o}gico real. Quando o programa resultante \'{e} executado na grade de teste, ele falha.

Para enfrentar esse problema, avaliamos neste trabalho uma estrat\'{e}gia com regras expl\'{\i}citas contra trapa\c{c}as l\'{o}gicas, denominada \textit{AntiCheat}. Essa estrat\'{e}gia imp\~{o}e duas restri\c{c}\~{o}es ao modelo no momento do feedback: (1) pro\'{\i}be o uso de condicionais direcionadas a coordenadas fixas ou \'{\i}ndices espec\'{\i}ficos de exemplos; e (2) exige que o modelo descreva em linguagem natural a regra abstrata geral antes de escrever a nova fun\c{c}\~{a}o em Python. 

A hip\'{o}tese central investigada neste trabalho \'{e}: \textit{adicionar regras que for\c{c}am explica\c{c}\~{o}es pr\'{e}vias e pro\'{\i}bem regras pontuais melhora a generaliza\c{c}\~{a}o e a acur\'{a}cia dos modelos no ARC-AGI?}

Compreender o impacto dessas restri\c{c}\~{o}es \'{e} fundamental para a engenharia de sistemas aut\^{o}nomos baseados em c\'{o}digo, pois ajuda a distinguir o racioc\'{\i}nio genu\'{\i}no da memoriza\c{c}\~{a}o superficial em modelos de linguagem.

\subsection*{Principais Contribui\c{c}\~{o}es}

As principais contribui\c{c}\~{o}es deste artigo s\~{a}o:
\begin{enumerate}
    \item \textbf{Avalia\c{c}\~{a}o estat\'{\i}stica rigorosa de CEGIS versus AntiCheat}: Realizamos 800 avalia\c{c}\~{o}es pareadas (8 modelos em 100 tarefas) utilizando testes de permuta\c{c}\~{a}o e o teste binomial exato para comparar diretamente a acur\'{a}cia das duas abordagens.
    \item \textbf{Inspe\c{c}\~{a}o qualitativa do c\'{o}digo gerado}: Analisamos as diferen\c{c}as reais no c\'{o}digo e os motivos de falha de cada abordagem, verificando se as restri\c{c}\~{o}es provocam mudan\c{c}as significativas na l\'{o}gica produzida.
    \item \textbf{Valida\c{c}\~{a}o da Escala de Guttman}: Demonstramos que o desempenho dos modelos no ARC-AGI forma uma escala psicom\'{e}trica unidimensional v\'{a}lida ($CR = 0,9738$), indicando que modelos mais fortes absorvem quase integralmente as capacidades dos modelos mais simples.
    \item \textbf{Comprova\c{c}\~{a}o emp\'{\i}rica da Navalha de Occam}: Mostramos, a partir da \'{a}rvore sint\'{a}tica abstrata (AST), que solu\c{c}\~{o}es com sobreajuste s\~{a}o 1,61 vezes mais extensas do que solu\c{c}\~{o}es corretas ($p = 5,15 \times 10^{-9}$).
    \item \textbf{Determina\c{c}\~{a}o do horizonte \'{o}timo de itera\c{c}\~{o}es}: Mapeamos a recupera\c{c}\~{a}o do CEGIS e mostramos que 86\% de todas as corre\c{c}\~{o}es acontecem at\'{e} a terceira itera\c{c}\~{a}o.
\end{enumerate}

\subsection*{Resultados Obtidos}

Adiantando os resultados, constatamos que a estrat\'{e}gia AntiCheat \textbf{n\~{a}o produz diferen\c{c}a estatisticamente significante} na acur\'{a}cia global em rela\c{c}\~{a}o ao CEGIS padr\~{a}o (50,4\% contra 49,8\%, com $p = 0,5515$). Al\'{e}m disso, a inspe\c{c}\~{a}o minuciosa do c\'{o}digo gerado revelou que, mesmo em n\'{\i}vel qualitativo, n\~{a}o h\'{a} diferen\c{c}a substancial: em mais de 88\% dos casos de discord\^{a}ncia, a estrat\'{e}gia perdedora falha simplesmente por n\~{a}o conseguir encontrar um algoritmo v\'{a}lido para os exemplos de treino, e n\~{a}o por tentar criar regras pontuais.

\subsection*{Estrutura do Artigo}

O restante do artigo est\'{a} organizado da seguinte forma: a Se\c{c}\~{a}o~\ref{sec:relacionados} aborda os trabalhos relacionados. A Se\c{c}\~{a}o~\ref{sec:proposta} descreve a abordagem do CEGIS e o mecanismo AntiCheat. A Se\c{c}\~{a}o~\ref{sec:metodologia} apresenta a metodologia experimental. A Se\c{c}\~{a}o~\ref{sec:resultados} detalha os resultados emp\'{\i}ricos e estruturais. Por fim, a Se\c{c}\~{a}o~\ref{sec:conclusao} discute as conclus\~{o}es e os trabalhos futuros.



% Seção 2: Trabalhos Relacionados
\section{Trabalhos Relacionados}
\label{sec:relacionados}

\subsection{O Benchmark ARC-AGI e a S\'{\i}ntese de Programas}

O \textit{Abstraction and Reasoning Corpus} (ARC) foi proposto por Chollet \cite{chollet2019measure} com o objetivo expl\'{\i}cito de medir a intelig\^{e}ncia geral artificial em contraste com a simples memoriza\c{c}\~{a}o de bases massivas de dados. No ARC, cada tarefa exige a indu\c{c}\~{a}o de conceitos como simetria, conectividade, trajet\'{o}ria e contagem a partir de 2 a 5 pares de treino. 

Trabalhos recentes mostram que a gera\c{c}\~{a}o direta de texto pixel a pixel por LLMs atinge desempenho limitado, ao passo que a s\'{\i}ntese de programas em linguagens como Python \cite{greenblatt2024getting} permite ao modelo delegar a execu\c{c}\~{a}o precisa e determin\'{\i}stica a um interpretador. Isso transforma o desafio do ARC em um problema cl\'{a}ssico de s\'{\i}ntese de programas a partir de exemplos.

\subsection{S\'{\i}ntese Indutiva Guiada por Contraexemplos (CEGIS)}

O paradigma CEGIS, formalizado por Solar-Lezama \cite{solar2008program}, estrutura a s\'{\i}ntese de software em um ciclo cooperativo entre dois componentes: um sintetizador e um verificador. O sintetizador gera um programa candidato, e o verificador busca entradas em que o programa falha (contraexemplos). Caso um contraexemplo seja encontrado, ele \'{e} adicionado \`{a} especifica\c{c}\~{a}o para que o sintetizador corrija o programa.

Com o advento dos LLMs, o CEGIS passou a ser empregado como um mecanismo interativo de corre\c{c}\~{a}o de c\'{o}digo \cite{poesia2024program}. Em vez de solucionadores formais (como SMT solvers), o pr\'{o}prio modelo atua como sintetizador probabil\'{\i}stico, recebendo mensagens de erro de execu\c{c}\~{a}o ou diverg\^{e}ncias de sa\'{\i}da para refinar sua solu\c{c}\~{a}o.

\subsection{Reflex\~{a}o e Racioc\'{\i}nio em LLMs}

Diversos m\'{e}todos prop\~{o}em que modelos de linguagem reflitam sobre seus pr\'{o}prios erros antes de gerar novas tentativas. Abordagens como \textit{Self-Refine} \cite{madaan2023self} e \textit{Reflexion} \cite{shinn2023reflexion} demonstraram que solicitar ao modelo um resumo verbal da falha pode aumentar a taxa de sucesso em problemas de racioc\'{\i}nio e programa\c{c}\~{a}o. A estrat\'{e}gia AntiCheat avaliada neste artigo segue essa linha, combinando a reflex\~{a}o verbal sobre invariantes com restri\c{c}\~{o}es expl\'{\i}citas contra regras pontuais.

\subsection{Escalabilidade Psicom\'{e}trica e Escala de Guttman}

A an\'{a}lise psicom\'{e}trica vem sendo progressivamente adotada para compreender a organiza\c{c}\~{a}o das habilidades de modelos de IA \cite{guttman1944basis}. A Escala de Guttman avalia se os itens de um teste seguem uma ordem cumulativa unidimensional de dificuldade: se uma tarefa \'{e} dominada por um modelo de habilidade m\'{e}dia, ela tamb\'{e}m deve ser dominada por modelos mais avan\c{c}ados. Aplicar essa an\'{a}lise ao ARC-AGI permite verificar se o benchmark reflete uma hierarquia de intelig\^{e}ncia consistente ou se o desempenho dos modelos \'{e} fragmentado e desordenado.



% Seção 3: Abordagem Proposta
\section{Abordagem Proposta}
\label{sec:proposta}

Nesta se\c{c}\~{a}o, descrevemos como formulamos o problema de s\'{\i}ntese de c\'{o}digo para o ARC-AGI, o funcionamento do ciclo CEGIS b\'{a}sico e a inclus\~{a}o das regras anti-ajuste da estrat\'{e}gia AntiCheat.

\subsection{Formula\c{c}\~{a}o do Problema}

Cada tarefa do ARC \'{e} composta por um conjunto de exemplos de treino $D_{\text{train}} = \{(I_1, O_1), \dots, (I_K, O_K)\}$ e um ou mais pares de teste $(I_{\text{test}}, O_{\text{test}})$, onde cada $I$ e $O$ s\~{a}o matrizes bidimensionais com n\'{u}meros inteiros no intervalo $[0, 9]$.

O objetivo \'{e} encontrar uma fun\c{c}\~{a}o em Python, denotada por $f(grid)$, que satisfa\c{c}a perfeitamente todos os exemplos de treino:
\begin{equation}
    f(I_k) = O_k, \quad \forall k \in \{1, \dots, K\}
\end{equation}
e que consiga generalizar para o exemplo de teste oculto:
\begin{equation}
    f(I_{\text{test}}) = O_{\text{test}}
\end{equation}

Para evitar o uso de bibliotecas externas complexas que poderiam mascarar o racioc\'{\i}nio do modelo, imp\~{o}e-se a regra de utilizar exclusivamente tipos e fun\c{c}\~{o}es nativas do Python (como listas, dicion\'{a}rios, conjuntos, la\c{c}os e fun\c{c}\~{o}es b\'{a}sicas como \texttt{len}, \texttt{range} e \texttt{min}).

\subsection{O Ciclo CEGIS Padr\~{a}o}

O algoritmo CEGIS opera em at\'{e} $T_{\max} = 5$ itera\c{c}\~{o}es. O processo funciona em quatro etapas simples:
\begin{enumerate}
    \item \textbf{Primeira Tentativa (\textit{Baseline})}: O modelo de linguagem recebe todos os pares de treino formatados em JSON e produz a primeira vers\~{a}o da fun\c{c}\~{a}o \texttt{transform(grid)}.
    \item \textbf{Execu\c{c}\~{a}o e Avalia\c{c}\~{a}o}: O c\'{o}digo gerado \'{e} executado de forma segura em cada exemplo $(I_k, O_k)$. Se o c\'{o}digo produzir a sa\'{\i}da correta para todos os $K$ exemplos de treino, o ciclo \'{e} interrompido com sucesso no treino (\textit{converg\^{e}ncia}).
    \item \textbf{Detec\c{c}\~{a}o do Contraexemplo}: Caso o c\'{o}digo erre em algum exemplo ou sofra um erro de execu\c{c}\~{a}o, o primeiro exemplo falho \'{e} selecionado como contraexemplo.
    \item \textbf{Feedback e Reparo}: Uma nova mensagem \'{e} enviada ao modelo contendo a entrada $I_k$, a sa\'{\i}da esperada $O_k$, a sa\'{\i}da incorretamente produzida (ou a mensagem de erro) e a instru\c{c}\~{a}o para corrigir o c\'{o}digo. O ciclo se repete at\'{e} atingir a converg\^{e}ncia ou esgotar as 5 itera\c{c}\~{o}es.
\end{enumerate}

\subsection{A Estrat\'{e}gia AntiCheat}

A estrat\'{e}gia AntiCheat introduz duas mudan\c{c}as fundamentais no prompt de feedback quando ocorre uma falha:

\begin{enumerate}
    \item \textbf{Proibi\c{c}\~{a}o de atalhos e regras pontuais}: A instru\c{c}\~{a}o pro\'{\i}be explicitamente que o modelo use desvios condicionais (\texttt{if/else}) voltados a \'{\i}ndices espec\'{\i}ficos de exemplos ou coordenadas fixas de uma grade isolada. O modelo \'{e} instru\'{\i}do a produzir uma regra matem\'{a}tica ou geom\'{e}trica uniforme que se aplique igualmente a todas as grades.
    \item \textbf{Enuncia\c{c}\~{a}o pr\'{e}via de invariantes}: Antes de gerar o bloco de c\'{o}digo Python, o modelo \'{e} obrigado a explicar em texto qual \'{e} o \textit{invariante abstrato} que falhou na vers\~{a}o anterior e como a nova l\'{o}gica atende a todos os exemplos conhecidos sem criar casos especiais.
\end{enumerate}

\subsection{Algoritmo Resumido}

O fluxo unificado de execu\c{c}\~{a}o \'{e} apresentado no Algoritmo~\ref{alg:cegis_anticheat}.

\begin{algorithm}[h]
\caption{Ciclo CEGIS com op\c{c}\~{a}o de regras AntiCheat}
\label{alg:cegis_anticheat}
\small
\textbf{Entrada:} Pares de treino $D_{\text{train}} = \{(I_k, O_k)\}_{k=1}^K$, Modo $\text{AntiCheat} \in \{\text{verdadeiro}, \text{falso}\}$, Limite $T_{\max} = 5$\\
\textbf{Sa\'{\i}da:} Programa final $f^*$ e status de converg\^{e}ncia

$f_1 \gets \text{GerarPrimeiraTentativa}(D_{\text{train}})$\;
\Para{$t \gets 1$ \Ate $T_{\max}$}{
    $\text{falhas} \gets \emptyset$\;
    \ParaCada{$(I_k, O_k) \in D_{\text{train}}$}{
        $\hat{O}_k \gets \text{Executar}(f_t, I_k)$\;
        \Se{$\hat{O}_k \neq O_k$}{
            $\text{falhas} \gets \text{falhas} \cup \{(k, I_k, O_k, \hat{O}_k)\}$\;
            \textbf{interromper}\; \tcp{Seleciona o primeiro contraexemplo}
        }
    }
    \Se{$\text{falhas} = \emptyset$}{
        \Retorna $f_t$, \text{convergiu} = verdadeiro\;
    }
    $(idx, I, O, \hat{O}) \gets \text{Primeiro}(\text{falhas})$\;
    $\text{prompt} \gets \text{MontarFeedback}(idx, I, O, \hat{O})$\;
    \Se{$\text{AntiCheat}$ \'{e} verdadeiro}{
        $\text{prompt} \gets \text{prompt} + \text{RegrasDeInvarianteEAntiAtalhos}()$\;
    }
    $f_{t+1} \gets \text{ChamarModelo}(\text{prompt})$\;
}
\Retorna $f_{T_{\max}}$, \text{convergiu} = falso\;
\end{algorithm}

\subsection{O Que H\'{a} de Novo}

A principal novidade desta proposta \'{e} o acoplamento entre um verificador determin\'{\i}stico estrito de c\'{o}digo (que n\~{a}o permite aproxima\c{c}\~{o}es) e restri\c{c}\~{o}es sem\^{a}nticas que impedem o modelo de trapacear na corre\c{c}\~{a}o de bugs. Em vez de permitir que o LLM resolva o erro por for\c{c}a bruta sint\'{a}tica, for\c{c}amos a gera\c{c}\~{a}o de uma hip\'{o}tese conceitual antes da escrita da l\'{o}gica em Python.



% Seção 4: Metodologia Experimental (deixada em branco para preenchimento posterior)
\section{Metodologia Experimental}
\label{sec:metodologia}

% [Seção deixada em branco conforme solicitado, a ser preenchida posteriormente.]



% Seção 5: Resultados
\section{Resultados}
\label{sec:resultados}

Nesta se\c{c}\~{a}o, apresentamos os resultados obtidos na avalia\c{c}\~{a}o de 8 modelos de linguagem em 100 tarefas do ARC-AGI, totalizando 800 pares de execu\c{c}\~{a}o. Analisamos a acur\'{a}cia das abordagens, o comportamento qualitativo do c\'{o}digo, a estrutura de dificuldade das tarefas, a complexidade do c\'{o}digo gerado e a din\^{a}mica do ciclo iterativo.

\subsection{Compara\c{c}\~{a}o de Acur\'{a}cia: CEGIS vs. AntiCheat}

A Tabela~\ref{tab:modelos} resume o desempenho de cada modelo sob as duas estrat\'{e}gias. No total de 800 avalia\c{c}\~{o}es, o CEGIS acertou 403 tarefas (50,4\%), enquanto o AntiCheat acertou 398 tarefas (49,8\%), representando uma diferen\c{c}a absoluta de apenas 5 tarefas ($-0,62$ pontos percentuais).

\begin{table}[h!]
\centering
\caption{Desempenho por modelo no CEGIS e no AntiCheat (100 tarefas por modelo).}
\label{tab:modelos}
\small
\begin{tabular}{lcccccc}
\toprule
\textbf{Modelo} & \textbf{CEGIS (\%)} & \textbf{AntiCheat (\%)} & \textbf{Dif. (\%)} & \textbf{Vit\'{o}rias ($b / c$)} & \textbf{$p$-valor bruto} & \textbf{$p$ ajustado (Bonferroni)} \\
\midrule
Grok 4.6 & 92,0\% & 92,0\% & 0,0\% & 0 / 0 & 1,0000 & 1,0000 \\
DeepSeek V4 Flash 0731 & 86,0\% & 86,0\% & 0,0\% & 0 / 0 & 1,0000 & 1,0000 \\
GPT-5.6 Luna & 77,0\% & 79,0\% & +2,0\% & 2 / 4 & 0,6875 & 1,0000 \\
Claude Sonnet 5 & 62,0\% & 60,0\% & -2,0\% & 5 / 3 & 0,7266 & 1,0000 \\
Gemini 3.5 Flash Lite & 43,0\% & 45,0\% & +2,0\% & 6 / 8 & 0,7905 & 1,0000 \\
Gemini 3.1 Flash Lite & 31,0\% & 24,0\% & -7,0\% & 7 / 0 & 0,0156 & 0,1250 \\
Leanstral 1.5.1 & 7,0\% & 9,0\% & +2,0\% & 2 / 4 & 0,6875 & 1,0000 \\
GPT-5.4 Nano & 5,0\% & 3,0\% & -2,0\% & 3 / 1 & 0,6250 & 1,0000 \\
\midrule
\textbf{Total Agregado} & \textbf{50,4\%} & \textbf{49,8\%} & \textbf{-0,6\%} & \textbf{25 / 20} & \textbf{0,5515} & --- \\
\bottomrule
\end{tabular}
\end{table}

Para verificar se essa diferen\c{c}a tem significado estat\'{\i}stico, aplicamos testes pareados nos pares discordantes (casos em que um m\'{e}todo acertou e o outro errou):
\begin{itemize}
    \item \textbf{Teste Binomial Exato}: Dos 45 casos de discord\^{a}ncia, o CEGIS venceu 25 e o AntiCheat venceu 20. O teste binomial exato resulta em $p = 0,5515$.
    \item \textbf{Teste de Permuta\c{c}\~{a}o}: Com 100.000 permuta\c{c}\~{o}es de sinal pareadas, obtivemos $p = 0,5501$.
    \item \textbf{Intervalo de Confian\c{c}a (\textit{Bootstrap})}: O intervalo de confian\c{c}a de 95\% com 10.000 reamostragens ficou entre $[-2,25\%, +1,00\%]$.
\end{itemize}

Como o valor de $p$ \'{e} substancialmente superior a $0,05$ e o intervalo de confian\c{c}a inclui o zero, \textbf{n\~{a}o h\'{a} evid\^{e}ncia estat\'{\i}stica de que a estrat\'{e}gia AntiCheat seja superior ou inferior ao CEGIS}.

Em rela\c{c}\~{a}o ao modelo \texttt{Gemini 3.1 Flash Lite}, que apresentou uma queda aparente de 7\% (7 acertos exclusivos do CEGIS contra 0 do AntiCheat, com $p = 0,0156$ n\~{a}o corrigido), a aplica\c{c}\~{a}o da corre\c{c}\~{a}o de Bonferroni para os 8 modelos avaliados eleva o $p$-valor para $0,1250$. Esse resultado neutraliza a diferen\c{c}a, indicando que a varia\c{c}\~{a}o observada \'{e} compat\'{\i}vel com a oscila\c{c}\~{a}o esperada do acaso ao testar m\'{u}ltiplas hip\'{o}teses simultaneamente, n\~{a}o se configurando como uma regress\~{a}o inesperada.

\subsection{An\'{a}lise Qualitativa do C\'{o}digo e Motivos de Falha}

Buscamos entender se, al\'{e}m dos n\'{u}meros de acur\'{a}cia, as regras do AntiCheat provocavam altera\c{c}\~{o}es qualitativas relevantes nos programas gerados. Inspecionamos detalhadamente os 45 casos discordantes e os motivos de falha de cada estrat\'{e}gia:

\begin{enumerate}
    \item \textbf{Motivo de falha predominante}: Nos 25 casos em que o CEGIS venceu, a falha do AntiCheat decorreu de teto de representa\c{c}\~{a}o (\textit{representation ceiling}) em 23 deles (92,0\%). Da mesma forma, nos 20 casos em que o AntiCheat venceu, o CEGIS falhou por teto de representa\c{c}\~{a}o em 18 deles (90,0\%). Em ambos os casos, a estrat\'{e}gia perdedora simplesmente n\~{a}o conseguiu encontrar uma l\'{o}gica que satisfizesse os exemplos de treino no limite de 5 itera\c{c}\~{o}es.
    \item \textbf{Ocorr\^{e}ncia de sobreajuste esp\'{u}rio}: O sobreajuste esp\'{u}rio (acertar o treino, mas errar o teste) ocorreu em apenas 2 tarefas de cada lado.
    \item \textbf{Aus\^{e}ncia de trapa\c{c}as por coordenadas}: Uma varredura est\'{a}tica em todas as 2.398 fun\c{c}\~{o}es geradas no benchmark revelou que apenas 7 c\'{o}digos (0,29\%) continham coordenadas fixas ou condicionais direcionadas a exemplos isolados.
\end{enumerate}

A inspe\c{c}\~{a}o do c\'{o}digo-fonte revelou que os dois m\'{e}todos produzem solu\c{c}\~{o}es estruturalmente muito pr\'{o}ximas. O modelo mant\'{e}m a mesma estrat\'{e}gia algor\'{\i}tmica b\'{a}sica em ambos os casos, e pequenas varia\c{c}\~{o}es no prompt afetam apenas detalhes pontuais de indexa\c{c}\~{a}o ou busca geom\'{e}trica. Portanto, \textbf{n\~{a}o se observa diferen\c{c}a qualitativa sistem\'{a}tica} entre o c\'{o}digo do CEGIS e o c\'{o}digo do AntiCheat.

\subsection{Dificuldade das Tarefas e Escala de Guttman}

A distribui\c{c}\~{a}o do n\'{u}mero de modelos que resolveram cada uma das 100 tarefas \'{e} apresentada na Figura~\ref{fig:task_dist}. Observamos que:
\begin{itemize}
    \item \textbf{Tarefas Universalmente Resolvidas}: Apenas 3 tarefas (3\%) foram resolvidas por todos os 8 modelos (\texttt{0c786b71}, \texttt{2072aba6}, \texttt{332efdb3}).
    \item \textbf{Tarefas Universalmente N\~{a}o Resolvidas}: 7 tarefas (7\%) n\~{a}o foram solucionadas por nenhum dos modelos avaliados.
    \item \textbf{Armadilhas Indutivas}: Foram identificadas 5 tarefas em que a maioria dos modelos atingiu 100\% de acerto nos treinos, mas falhou totalmente no teste (por exemplo, a tarefa \texttt{351d6448}, onde 6 dos 8 modelos convergiram no treino, mas nenhum acertou o teste).
\end{itemize}

\begin{figure}[h!]
\centering
\includegraphics[width=0.75\linewidth]{static_analysis_output/plot_task_solve_distribution.png}
\caption{Distribui\c{c}\~{a}o de tarefas resolvidas por quantidade de modelos.}
\label{fig:task_dist}
\end{figure}

Para analisar se o benchmark mede uma habilidade cumulativa \'{u}nica, calculamos a Matriz de Subsun\c{c}\~{a}o de Guttman, ilustrada na Figura~\ref{fig:guttman}. A probabilidade de um modelo superior acertar as tarefas que um modelo inferior acerta \'{e} expressiva:
\begin{itemize}
    \item \texttt{Grok 4.6} resolveu \textbf{100,0\%} das tarefas resolvidas pelo \texttt{DeepSeek V4 Flash};
    \item \texttt{DeepSeek V4 Flash} resolveu \textbf{97,4\%} das tarefas resolvidas pelo \texttt{GPT-5.6 Luna};
    \item \texttt{GPT-5.6 Luna} resolveu \textbf{91,9\%} das tarefas resolvidas pelo \texttt{Claude Sonnet 5};
    \item \texttt{Claude Sonnet 5} resolveu \textbf{86,0\%} das tarefas resolvidas pelo \texttt{Gemini 3.5 Flash Lite}.
\end{itemize}

\begin{figure}[h!]
\centering
\includegraphics[width=0.75\linewidth]{static_analysis_output/plot_guttman_subsumption.png}
\caption{Matriz de Subsun\c{c}\~{a}o de Guttman entre os 8 modelos avaliados.}
\label{fig:guttman}
\end{figure}

Os \'{\i}ndices psicom\'{e}tricos calculados confirmam que o ARC-AGI forma uma escala unidimensional quase perfeita:
\begin{itemize}
    \item \textbf{Coeficiente de Reprodutibilidade (CR)}: $0,9738$ (o limiar de aceita\c{c}\~{a}o psicom\'{e}trico \'{e} $\ge 0,90$);
    \item \textbf{Coeficiente de Escalabilidade (CS)}: $0,8757$ (o limiar de aceita\c{c}\~{a}o \'{e} $\ge 0,60$).
\end{itemize}
Isso confirma que o ganho de desempenho dos modelos segue uma progress\~{a}o estritamente hier\'{a}rquica de compet\^{e}ncias.

\subsection{Complexidade do C\'{o}digo e a Navalha de Occam}

Analisamos a \'{A}rvore Sint\'{a}tica Abstrata (AST) dos programas gerados para investigar a rela\c{c}\~{a}o entre o tamanho do c\'{o}digo e o sucesso da solu\c{c}\~{a}o (Figura~\ref{fig:complexity}).

\begin{table}[h!]
\centering
\caption{Complexidade estrutural do c\'{o}digo por desfecho da execu\c{c}\~{a}o.}
\label{tab:complexity}
\small
\begin{tabular}{lccccc}
\toprule
\textbf{Categoria de Desfecho} & \textbf{Amostras} & \textbf{Linhas (SLOC)} & \textbf{La\c{c}os (\textit{loops})} & \textbf{Condicionais (\textit{ifs})} & \textbf{Profundidade AST} \\
\midrule
Acerto de Primeira (\textit{both correct}) & 289 & 34,6 & 6,2 & 4,9 & 12,2 \\
Recuperado via CEGIS (\textit{semantic recovery}) & 114 & 44,5 & 7,6 & 7,0 & 12,7 \\
Sobreajuste Esp\'{u}rio (\textit{spurious overfitting}) & 52 & 55,9 & 8,8 & 8,7 & 13,1 \\
Falha no Treino (\textit{ceiling}) & 342 & 44,4 & 7,9 & 7,5 & 12,8 \\
\bottomrule
\end{tabular}
\end{table}

\begin{figure}[h!]
\centering
\includegraphics[width=0.75\linewidth]{static_analysis_output/plot_code_complexity_by_outcome.png}
\caption{Distribui\c{c}\~{a}o das linhas de c\'{o}digo (SLOC) de acordo com o desfecho da tarefa.}
\label{fig:complexity}
\end{figure}

Como resumido na Tabela~\ref{tab:complexity}, os c\'{o}digos que acertam na primeira tentativa possuem m\'{e}dia de 34,6 linhas de c\'{o}digo. Em contraste, os c\'{o}digos que caem em sobreajuste esp\'{u}rio apresentam uma m\'{e}dia de 55,9 linhas. Esse aumento representa um fator de incha\c{c}o (\textit{Bloat Factor}) de \textbf{1,61 vezes} ($p = 5,15 \times 10^{-9}$ no teste de Mann-Whitney U). Esse resultado fornece uma valida\c{c}\~{a}o emp\'{\i}rica rigorosa do princ\'{\i}pio da \textbf{Navalha de Occam}: quando os modelos tentam acomodar contraexemplos sem entender o padr\~{a}o geral, tendem a inflar o programa com condicionais e la\c{c}os adicionais que n\~{a}o generalizam.

\subsection{Din\^{a}mica de Itera\c{c}\~{o}es do CEGIS}

Avaliamos quantas itera\c{c}\~{o}es s\~{a}o necess\'{a}rias para que o CEGIS conserte c\'{o}digos que falharam na primeira tentativa. Conforme apresentado na Tabela~\ref{tab:iterations} e na Figura~\ref{fig:iterations}:
\begin{itemize}
    \item 71,7\% (289 tarefas) do total de acertos ocorrem na primeira itera\c{c}\~{a}o;
    \item Para as 114 tarefas corrigidas ao longo do ciclo iterativo, 60,5\% das corre\c{c}\~{o}es acontecem logo na 2\textordfeminine{} itera\c{c}\~{a}o e 25,4\% ocorrem na 3\textordfeminine{} itera\c{c}\~{a}o;
    \item As itera\c{c}\~{o}es 4 e 5 respondem por apenas 7,9\% e 6,1\% das corre\c{c}\~{o}es, respectivamente.
\end{itemize}

\begin{table}[h!]
\centering
\caption{Distribui\c{c}\~{a}o de recupera\c{c}\~{a}o de tarefas ao longo das itera\c{c}\~{o}es do CEGIS.}
\label{tab:iterations}
\small
\begin{tabular}{ccccc}
\toprule
\textbf{Itera\c{c}\~{a}o} & \textbf{Tarefas Resolvidas} & \textbf{\% do Total Resolvido} & \textbf{Tarefas Recuperadas} & \textbf{\% Acumulado de Recupera\c{c}\~{o}es} \\
\midrule
1 & 289 & 71,7\% & 0 & 0,0\% \\
2 & 69 & 17,1\% & 69 & 60,5\% \\
3 & 29 & 7,2\% & 29 & \textbf{86,0\%} \\
4 & 9 & 2,2\% & 9 & 93,9\% \\
5 & 7 & 1,7\% & 7 & 100,0\% \\
\bottomrule
\end{tabular}
\end{table}

\begin{figure}[h!]
\centering
\includegraphics[width=0.75\linewidth]{static_analysis_output/plot_iteration_diminishing_returns.png}
\caption{Rendimentos decrescentes ao longo das itera\c{c}\~{o}es do CEGIS.}
\label{fig:iterations}
\end{figure}

Esse comportamento demonstra um retorno decrescente r\'{a}pido. Fixar o limite do CEGIS em \textbf{3 itera\c{c}\~{o}es} permite capturar 86,0\% de todos os reparos poss\'{\i}veis, evitando o consumo desnecess\'{a}rio de tokens e processamento nas rodadas finais.



% Seção 6: Discussão e Conclusões
\section{Discuss\~{a}o e Conclus\~{o}es}
\label{sec:conclusao}

\subsection{S\'{\i}ntese dos Resultados e Interpreta\c{c}\~{a}o}

Neste trabalho, investigamos se a introdu\c{c}\~{a}o de restri\c{c}\~{o}es sem\^{a}nticas expl\'{\i}citas contra regras pontuais e a obrigatoriedade de explicar invariantes (\textit{AntiCheat}) trariam ganhos reais sobre a s\'{\i}ntese iterativa de programas por CEGIS no benchmark ARC-AGI.

A principal conclus\~{a}o emp\'{\i}rica \'{e} que a estrat\'{e}gia AntiCheat \textbf{n\~{a}o produziu diferen\c{c}a estatisticamente significante} na acur\'{a}cia geral (50,4\% contra 49,8\%, com $p = 0,5515$). Al\'{e}m disso, a inspe\c{c}\~{a}o qualitativa confirmou que essa equival\^{e}ncia ocorre porque as falhas em ambos os m\'{e}todos s\~{a}o causadas quase que exclusivamente pelo mesmo motivo: o esgotamento da capacidade do modelo de formular o algoritmo correto (\textit{representation ceiling}), e n\~{a}o pela tentativa deliberada de criar atalhos com coordenadas fixas. Em modelos avan\c{c}ados de linguagem, a cria\c{c}\~{a}o de atalhos pontuais de coordenadas \'{e} um fen\^{o}meno residual (presente em menos de 0,3\% dos c\'{o}digos analisados).

Por outro lado, nossa an\'{a}lise estrutural produziu descobertas relevantes e independentes de estrat\'{e}gias espec\'{\i}ficas:
\begin{enumerate}
    \item \textbf{Consist\^{e}ncia Psicom\'{e}trica}: O ARC-AGI exibe uma organiza\c{c}\~{a}o hier\'{a}rquica de dificuldade quase perfeita na Escala de Guttman ($CR = 0,9738$), mostrando que o benchmark reflete uma dimens\~{a}o est\'{a}vel de avan\c{c}o de capacidades dos modelos.
    \item \textbf{Valida\c{c}\~{a}o da Navalha de Occam}: Programas que caem em sobreajuste sofrem uma infla\c{c}\~{a}o sistem\'{a}tica no n\'{u}mero de linhas (1,61 vezes maiores, $p = 5,15 \times 10^{-9}$), la\c{c}os e desvios condicionais, confirmando numericamente que solu\c{c}\~{o}es simples t\^{e}m maior probabilidade de generaliza\c{c}\~{a}o.
    \item \textbf{Horizonte de Efici\^{e}ncia}: Estabelecemos que limitar o ciclo de corre\c{c}\~{a}o do CEGIS a 3 itera\c{c}\~{o}es preserva 86,0\% dos reparos vi\'{a}veis, otimizando o custo computacional.
\end{enumerate}

\subsection{Aplica\c{c}\~{o}es e Impactos}

Estes achados trazem orienta\c{c}\~{o}es diretas para o desenvolvimento de sistemas de programa\c{c}\~{a}o assistida por IA:
\begin{itemize}
    \item \textbf{Crit\'{e}rios de sele\c{c}\~{a}o por simplicidade}: Sistemas de busca e gera\c{c}\~{a}o de c\'{o}digo devem usar o tamanho da AST e o n\'{u}mero de linhas como fun\c{c}\~{a}o de penalidade para descartar programas que se tornam excessivamente complexos ap\'{o}s m\'{u}ltiplas tentativas de corre\c{c}\~{a}o.
    \item \textbf{Crit\'{e}rios de parada precoce}: Encerrar o ciclo CEGIS ap\'{o}s 3 itera\c{c}\~{o}es \'{e} uma pol\'{\i}tica eficiente para evitar gasto de recursos em solu\c{c}\~{o}es que raramente convergem.
\end{itemize}

\subsection{Limita\c{c}\~{o}es}

As principais limita\c{c}\~{o}es deste estudo incluem:
\begin{enumerate}
    \item \textbf{Amostragem de tarefas}: A avalia\c{c}\~{a}o utilizou uma amostra representativa de 100 tarefas do conjunto de treino/avalia\c{c}\~{a}o do ARC-AGI.
    \item \textbf{Dom\'{\i}nio restrito ao Python nativo}: A regra de n\~{a}o permitir bibliotecas externas padroniza o teste, mas pode restringir o estilo de solu\c{c}\~{a}o dos modelos.
    \item \textbf{Temperatura e amostragem \'{u}nica}: Cada tarefa foi executada com par\^{a}metros fixos de amostragem, sem explora\c{c}\~{a}o massiva de sementes aleat\'{o}rias (\textit{pass@k}).
\end{enumerate}

\subsection{Trabalhos Futuros}

Como perspectivas de trabalhos futuros, sugerimos:
\begin{itemize}
    \item \textbf{Varia\c{c}\~{a}o sistem\'{a}tica entre prompts}: Investigar como diferentes estilos de instru\c{c}\~{a}o, exemplos no prompt (\textit{few-shot prompting}) e formula\c{c}\~{o}es alternativas de regras afetam a capacidade do modelo de se autorregular.
    \item \textbf{Pol\'{\i}ticas adaptativas de corre\c{c}\~{a}o}: Desenvolver mecanismos que decidem o tipo de feedback com base na complexidade do erro observado na execu\c{c}\~{a}o anterior.
    \item \textbf{Filtros de sele\c{c}\~{a}o sint\'{a}tica}: Aplicar limites r\'{\i}gidos de linhas de c\'{o}digo durante a s\'{\i}ntese para barrar automaticamente c\'{o}digos que sofram o incha\c{c}o caracter\'{\i}stico do sobreajuste.
\end{itemize}



% Seção 7: Referências Bibliográficas
\bibliography{references}

\end{document}
